% =============================================================================
% Aula 01 — Introdução ao Spring Boot + Setup + Primeiro projeto
% Beamer + tema Madrid | Curso Entra21 — Spring Framework
% Compilar: pdflatex slides-aula-01.tex  (duas vezes se usar TOC)
% =============================================================================
\documentclass[aspectratio=169,11pt]{beamer}

\usetheme{Madrid}
\usecolortheme{default}

\usepackage[T1]{fontenc}
\usepackage[utf8]{inputenc}
\usepackage[brazil]{babel}
\usepackage{booktabs}
\usepackage{graphicx}
\usepackage{hyperref}
\usepackage{listings}
\usepackage{xcolor}
\usepackage{tikz}
\usetikzlibrary{arrows.meta, positioning, shapes.geometric}

% --- Cores e código ---
\definecolor{codebg}{RGB}{245,245,245}
\definecolor{codegreen}{RGB}{40,120,40}
\definecolor{codeblue}{RGB}{30,70,160}
\definecolor{codered}{RGB}{180,50,50}

\lstset{
  basicstyle=\ttfamily\small,
  backgroundcolor=\color{codebg},
  frame=single,
  rulecolor=\color{black!20},
  breaklines=true,
  columns=fullflexible,
  keepspaces=true,
  showstringspaces=false,
  keywordstyle=\color{codeblue}\bfseries,
  commentstyle=\color{codegreen},
  stringstyle=\color{codered},
  tabsize=2,
  literate={á}{{\'a}}1 {é}{{\'e}}1 {í}{{\'i}}1 {ó}{{\'o}}1 {ú}{{\'u}}1
           {Á}{{\'A}}1 {É}{{\'E}}1 {Í}{{\'I}}1 {Ó}{{\'O}}1 {Ú}{{\'U}}1
           {ã}{{\~a}}1 {õ}{{\~o}}1 {Ã}{{\~A}}1 {Õ}{{\~O}}1
           {ç}{{\c{c}}}1 {Ç}{{\c{C}}}1
}

\title{Aula 01 --- Spring Framework}
\subtitle{Introdução · Setup no Windows · Primeiro projeto}
\author{Entra21 --- Trilha Java}
\institute{Última disciplina da trilha}
\date{}

\begin{document}

% =============================================================================
\begin{frame}
  \titlepage
\end{frame}

% =============================================================================
\begin{frame}{Bem-vindos à última disciplina da trilha}
  Vocês já passaram por muita coisa:

  \vspace{0.5em}
  \begin{enumerate}
    \item Fundamentos de Programação
    \item Orientação a Objetos
    \item Java Web
    \item \textbf{Spring Framework} $\leftarrow$ estamos aqui
  \end{enumerate}

  \vspace{0.8em}
  \begin{block}{Promessa desta disciplina}
    Vamos criar aplicações web em Java de um jeito \textbf{mais prático}
    do que montar Servlet + Tomcat na mão --- sem apagar o que vocês
    já aprenderam.
  \end{block}
\end{frame}

% =============================================================================
\begin{frame}{Como vamos estudar}
  \begin{itemize}
    \item Ritmo \textbf{devagar} e com pausas para dúvidas
    \item Explicações com \textbf{analogias} do dia a dia
    \item Um único repositório com \textbf{todos os códigos} da disciplina
    \item Em cada aula: algo feito \textbf{juntos} + um \textbf{desafio} sozinho
  \end{itemize}

  \vspace{0.8em}
  \begin{alertblock}{Regra de ouro}
    Se travou: pare, respire, peça ajuda. Travamento faz parte do processo.
  \end{alertblock}
\end{frame}

% =============================================================================
\begin{frame}{Mapa da aula de hoje}
  \tableofcontents
\end{frame}

% #############################################################################
\section{O que é Spring Boot?}

% =============================================================================
\begin{frame}{Lembra do Java Web?}
  Na disciplina anterior, para uma página responder no navegador, vocês
  tipicamente precisavam cuidar de várias peças:

  \vspace{0.6em}
  \begin{itemize}
    \item Instalar / configurar o \textbf{Tomcat}
    \item Criar \textbf{Servlets}
    \item Configurar o \texttt{web.xml} (em muitos projetos)
    \item Empacotar e publicar a aplicação
  \end{itemize}

  \vspace{0.6em}
  Funciona. Mas tem muita \textbf{montagem} antes do “olá, mundo”.
\end{frame}

% =============================================================================
\begin{frame}{Analogia da cozinha}
  \begin{columns}[T]
    \begin{column}{0.48\textwidth}
      \begin{block}{Java Web ``na mão''}
        Você monta a cozinha:\\
        fogão, pia, geladeira\ldots\\
        e \textbf{depois} cozinha.
      \end{block}
    \end{column}
    \begin{column}{0.48\textwidth}
      \begin{block}{Spring Boot}
        A cozinha já vem montada.\\
        Você foca em \textbf{preparar o prato}\\
        (a regra de negócio).
      \end{block}
    \end{column}
  \end{columns}

  \vspace{1em}
  \begin{alertblock}{Frase para lembrar}
    Spring Boot \textbf{não substitui} Java.\\
    Ele \textbf{organiza e acelera} o trabalho com Java.
  \end{alertblock}
\end{frame}

% =============================================================================
\begin{frame}{Spring $\times$ Spring Boot}
  \begin{description}
    \item[Spring] Conjunto grande de ferramentas para aplicações Java
      (web, dados, segurança\ldots).
    \item[Spring Boot] Forma \textbf{prática} de usar o Spring:
      configurações inteligentes + servidor embutido + projeto pronto para rodar.
  \end{description}

  \vspace{0.8em}
  Nesta disciplina, vamos falar principalmente de \textbf{Spring Boot}.\\
  Quando alguém no mercado diz “trabalhei com Spring”, quase sempre
  envolve Spring Boot.
\end{frame}

% =============================================================================
\begin{frame}{O que o Spring Boot faz por nós (versão simples)}
  \begin{enumerate}
    \item Sobe um \textbf{servidor web embutido} (não precisamos instalar Tomcat separado)
    \item Lê o \texttt{pom.xml} e baixa as \textbf{dependências}
    \item Liga as peças da aplicação com menos arquivo de configuração manual
    \item Deixa a gente escrever classes com anotações como \texttt{@RestController}
  \end{enumerate}

  \vspace{0.8em}
  \begin{block}{Meta de hoje}
    Criar um projeto, rodar, e ver uma mensagem no navegador.\\
    Só isso --- e isso já é uma vitória grande.
  \end{block}
\end{frame}

% =============================================================================
\begin{frame}{Comparação rápida (mental)}
  \begin{center}
  \begin{tabular}{lll}
    \toprule
    Ideia & Java Web & Spring Boot \\
    \midrule
    Pedido HTTP & Servlet & Controller \\
    Servidor & Tomcat externo & Embutido \\
    Configuração & muita manual & muita automática \\
    Dependências & Maven (\texttt{pom.xml}) & Maven (\texttt{pom.xml}) \\
    \bottomrule
  \end{tabular}
  \end{center}

  \vspace{1em}
  \textbf{Maven continua!} Vocês já viram \texttt{pom.xml}. Isso ajuda bastante.
\end{frame}

% #############################################################################
\section{Setup no Windows}

% =============================================================================
\begin{frame}{Checklist de instalação (Windows)}
  Vamos instalar / conferir \textbf{três coisas}:

  \vspace{0.6em}
  \begin{enumerate}
    \item \textbf{JDK 21} --- a máquina virtual Java
    \item \textbf{VS Code} --- nosso editor
    \item Extensões: \textbf{Java} + \textbf{Spring Boot}
  \end{enumerate}

  \vspace{0.8em}
  \begin{alertblock}{Importante}
    Façam \textbf{juntos} com o professor. Não adianta correr e pular etapa.\\
    Cada máquina é um pouco diferente --- vamos com calma.
  \end{alertblock}
\end{frame}

% =============================================================================
\begin{frame}{Passo 1 --- Instalar o JDK 21}
  \begin{enumerate}
    \item Abra o navegador em:\\
          \url{https://adoptium.net/}
    \item Escolha:
      \begin{itemize}
        \item Version: \textbf{21 (LTS)}
        \item Operating System: \textbf{Windows}
        \item Architecture: \textbf{x64} (maioria dos PCs)
      \end{itemize}
    \item Baixe o instalador \texttt{.msi} e execute
    \item Aceite as opções padrão (deixe marcar \textbf{Add to PATH} se aparecer)
  \end{enumerate}
\end{frame}

% =============================================================================
\begin{frame}[fragile]{Passo 1 --- Conferir o Java}
  Abra o \textbf{Prompt de Comando} ou o \textbf{Terminal} do Windows:

  \vspace{0.4em}
\begin{lstlisting}[language=bash]
java -version
\end{lstlisting}

  \vspace{0.4em}
  Você deve ver algo parecido com:

\begin{lstlisting}[basicstyle=\ttfamily\scriptsize]
openjdk version "21.0.x" ...
\end{lstlisting}

  \vspace{0.5em}
  \begin{block}{Se der ``não é reconhecido como comando''}
    O Java não entrou no PATH. Reinstale marcando a opção de PATH\\
    ou peça ajuda ao professor antes de continuar.
  \end{block}
\end{frame}

% =============================================================================
\begin{frame}{Passo 2 --- Instalar o VS Code}
  \begin{enumerate}
    \item Acesse \url{https://code.visualstudio.com/}
    \item Baixe para \textbf{Windows}
    \item Instale com as opções padrão
    \item (Opcional, mas útil) marque:
      \begin{itemize}
        \item Abrir com Code no menu de contexto
        \item Adicionar ao PATH
      \end{itemize}
  \end{enumerate}

  \vspace{0.6em}
  O VS Code é só o editor. As “superpoderes” Java vêm das \textbf{extensões}.
\end{frame}

% =============================================================================
\begin{frame}{Passo 3 --- Extensões no VS Code}
  No VS Code, clique no ícone de \textbf{Extensões} (ou \texttt{Ctrl+Shift+X}).

  \vspace{0.5em}
  Instale estas duas:

  \begin{enumerate}
    \item \textbf{Extension Pack for Java} (Microsoft)
    \item \textbf{Spring Boot Extension Pack} (VMware / Spring)
  \end{enumerate}

  \vspace{0.6em}
  Espere terminar o download (pode demorar um pouco na primeira vez).

  \vspace{0.4em}
  \begin{alertblock}{Dica}
    Reinicie o VS Code depois de instalar as extensões.
  \end{alertblock}
\end{frame}

% =============================================================================
\begin{frame}{Preciso instalar Maven separado?}
  \textbf{Não necessariamente.}

  \vspace{0.5em}
  Os projetos Spring Boot vêm com o \textbf{Maven Wrapper}:

  \begin{itemize}
    \item No Windows: \texttt{mvnw.cmd}
    \item Ele baixa o Maven automaticamente na primeira execução
  \end{itemize}

  \vspace{0.6em}
  \begin{block}{Tradução}
    Vocês digitam um comando na pasta do projeto.\\
    O próprio projeto cuida do Maven.
  \end{block}
\end{frame}

% =============================================================================
\begin{frame}{Checkpoint --- estamos prontos?}
  Antes do primeiro projeto, confirme:

  \vspace{0.5em}
  \begin{itemize}
    \item [ ] \texttt{java -version} mostra 21 (ou 17+)
    \item [ ] VS Code abre normalmente
    \item [ ] Extension Pack for Java instalado
    \item [ ] Spring Boot Extension Pack instalado
  \end{itemize}

  \vspace{0.8em}
  Se algum item falhou: \textbf{chame o professor agora}.\\
  Não siga adiante “no escuro”.
\end{frame}

% #############################################################################
\section{Primeiro projeto (juntos)}

% =============================================================================
\begin{frame}{O que vamos construir juntos}
  Um projeto chamado \textbf{ola-spring} que:

  \vspace{0.5em}
  \begin{enumerate}
    \item Sobe um servidor na porta \textbf{8080}
    \item Responde em \texttt{http://localhost:8080/}
    \item Responde também em \texttt{http://localhost:8080/ola}
  \end{enumerate}

  \vspace{0.8em}
  \begin{block}{Sucesso de hoje}
    Ver uma frase escrita por vocês no navegador.\\
    Isso prova que o Spring Boot está vivo na máquina.
  \end{block}
\end{frame}

% =============================================================================
\begin{frame}{Passo A --- Abrir o Spring Initializr}
  O site que \textbf{gera o esqueleto} do projeto:

  \vspace{0.4em}
  \begin{center}
    \Large \url{https://start.spring.io}
  \end{center}

  \vspace{0.8em}
  Pensem nele como um “formulário de pedido”:\\
  vocês escolhem o que querem, e ele monta a pasta inicial.
\end{frame}

% =============================================================================
\begin{frame}{Passo B --- Preencher o formulário}
  Configure assim (pode ir lendo em voz alta com a turma):

  \vspace{0.4em}
  \begin{itemize}
    \item \textbf{Project:} Maven
    \item \textbf{Language:} Java
    \item \textbf{Spring Boot:} deixe a versão padrão do site
    \item \textbf{Group:} \texttt{com.entra21}
    \item \textbf{Artifact:} \texttt{ola-spring}
    \item \textbf{Name:} \texttt{ola-spring}
    \item \textbf{Package name:} \texttt{com.entra21.olaspring}
    \item \textbf{Packaging:} Jar
    \item \textbf{Java:} 21
  \end{itemize}
\end{frame}

% =============================================================================
\begin{frame}{Passo C --- Adicionar a dependência certa}
  Em \textbf{Dependencies}, clique em \textbf{ADD DEPENDENCIES} e busque:

  \vspace{0.5em}
  \begin{center}
    \Large \textbf{Spring Web}
  \end{center}

  \vspace{0.5em}
  Adicione \textbf{só essa} por enquanto.

  \vspace{0.6em}
  \begin{alertblock}{Por quê?}
    Sem Spring Web, o projeto nasce --- mas \textbf{não} sobe endpoint HTTP
    pronto para o navegador do jeito que queremos hoje.
  \end{alertblock}
\end{frame}

% =============================================================================
\begin{frame}{Passo D --- Gerar e abrir no VS Code}
  \begin{enumerate}
    \item Clique em \textbf{GENERATE}
    \item Baixe o arquivo \texttt{.zip}
    \item Extraia para uma pasta fácil, por exemplo:\\
          \texttt{C:\textbackslash{}projetos\textbackslash{}ola-spring}
    \item No VS Code: \textbf{File $\rightarrow$ Open Folder}
    \item Selecione a pasta \texttt{ola-spring} (a que contém o \texttt{pom.xml})
  \end{enumerate}

  \vspace{0.5em}
  \begin{block}{Armadilha comum}
    Abrir a pasta \emph{pai} em vez da pasta do projeto.\\
    O VS Code precisa “enxergar” o \texttt{pom.xml} na raiz aberta.
  \end{block}
\end{frame}

% =============================================================================
\begin{frame}[fragile]{Passo E --- Olhar a estrutura (sem medo)}
\begin{lstlisting}[basicstyle=\ttfamily\scriptsize]
ola-spring/
├── pom.xml                 ← dependências (Maven)
├── mvnw.cmd                ← Maven Wrapper (Windows)
└── src/main/java/com/entra21/olaspring/
    └── OlaSpringApplication.java   ← ponto de partida (main)
\end{lstlisting}

  \vspace{0.5em}
  Hoje só vamos \textbf{criar mais uma classe}: o Controller.
\end{frame}

% =============================================================================
\begin{frame}[fragile]{Passo F --- Criar o OlaController}
  Na mesma pasta da Application, crie \texttt{OlaController.java}:

\begin{lstlisting}[language=Java,basicstyle=\ttfamily\footnotesize]
package com.entra21.olaspring;

import org.springframework.web.bind.annotation.GetMapping;
import org.springframework.web.bind.annotation.RestController;

@RestController
public class OlaController {

  @GetMapping("/")
  public String home() {
    return "Ola, Spring Boot! Bem-vindo(a).";
  }

  @GetMapping("/ola")
  public String ola() {
    return "Tudo funcionando!";
  }
}
\end{lstlisting}
\end{frame}

% =============================================================================
\begin{frame}{O que essas anotações significam?}
  \begin{description}
    \item[\texttt{@RestController}]
      “Esta classe recebe pedidos HTTP e devolve resposta.”
    \item[\texttt{@GetMapping("/ola")}]
      “Quando alguém acessar \texttt{/ola} com GET, rode este método.”
  \end{description}

  \vspace{0.8em}
  \begin{block}{Ponte com Java Web}
    Isso lembra um Servlet mapeado para uma URL ---\\
    só que com bem menos “cerimônia”.
  \end{block}
\end{frame}

% =============================================================================
\begin{frame}[fragile]{Passo G --- Rodar o projeto}
  No terminal do VS Code (\texttt{Ctrl+''} ou menu Terminal):

\begin{lstlisting}[language=bash]
.\mvnw.cmd spring-boot:run
\end{lstlisting}

  \vspace{0.4em}
  \begin{itemize}
    \item A \textbf{primeira} vez demora (baixa bibliotecas)
    \item Espere aparecer algo como:\\
          \texttt{Started OlaSpringApplication}
  \end{itemize}

  \vspace{0.5em}
  \begin{alertblock}{Não feche o terminal}
    Enquanto o servidor estiver rodando, o terminal fica “ocupado”.\\
    Isso é normal. Para parar: \texttt{Ctrl+C}.
  \end{alertblock}
\end{frame}

% =============================================================================
\begin{frame}{Passo H --- Testar no navegador}
  Com o servidor rodando, abra:

  \vspace{0.5em}
  \begin{itemize}
    \item \url{http://localhost:8080/}
    \item \url{http://localhost:8080/ola}
  \end{itemize}

  \vspace{0.8em}
  \begin{block}{O que é localhost:8080?}
    \begin{itemize}
      \item \textbf{localhost} = este computador
      \item \textbf{8080} = porta onde o Spring está escutando
    \end{itemize}
  \end{block}
\end{frame}

% =============================================================================
\begin{frame}{Se algo der errado\ldots}
  \begin{enumerate}
    \item Leia a mensagem vermelha do terminal (de baixo para cima)
    \item Confira se o pacote das classes está igual (\texttt{com.entra21.olaspring})
    \item Confira se a dependência \textbf{Spring Web} foi adicionada
    \item Veja se a porta 8080 já está em uso (feche outro Tomcat/Spring)
    \item Peça ajuda --- mostre a tela do erro
  \end{enumerate}

  \vspace{0.6em}
  Erro não é fracasso. Erro é o computador pedindo ajuste.
\end{frame}

% =============================================================================
\begin{frame}{Fluxo mental do que aconteceu}
  \begin{center}
  \begin{tikzpicture}[
      box/.style={draw, rounded corners, align=center, minimum width=2.6cm,
                  minimum height=1.0cm, font=\small\bfseries},
      seta/.style={-{Latex}, thick}
    ]
    \node[box, fill=blue!15] (nav) {Navegador\\{\normalfont\tiny localhost:8080}};
    \node[box, fill=orange!20, right=1.1cm of nav] (boot) {Spring Boot\\{\normalfont\tiny servidor}};
    \node[box, fill=green!15, right=1.1cm of boot] (ctrl) {OlaController\\{\normalfont\tiny @GetMapping}};
    \node[box, fill=yellow!25, below=0.9cm of ctrl] (txt) {Texto\\{\normalfont\tiny resposta}};

    \draw[seta] (nav) -- (boot);
    \draw[seta] (boot) -- (ctrl);
    \draw[seta] (ctrl) -- (txt);
    \draw[seta] (txt.west) -| ++(-1.2,0) |- (nav.south);
  \end{tikzpicture}
  \end{center}

  \vspace{0.6em}
  Vocês pedem $\rightarrow$ Spring recebe $\rightarrow$ Controller responde $\rightarrow$ navegador mostra.
\end{frame}

% #############################################################################
\section{Exercícios na sala}

% =============================================================================
\begin{frame}{Hora de praticar --- no ola-spring}
  Vamos treinar \textbf{no mesmo projeto} \texttt{ola-spring},\\
  criando novos métodos no \texttt{OlaController}.

  \vspace{0.6em}
  \begin{block}{Regras de ouro}
    \begin{enumerate}
      \item Faça \textbf{um exercício por vez}
      \item Salve o arquivo
      \item \textbf{Pare e rode de novo} o Spring (play ou terminal)
      \item Teste no navegador: \texttt{http://localhost:8080/...}
    \end{enumerate}
  \end{block}

  \vspace{0.4em}
  \begin{alertblock}{Lembrete}
    Se a resposta no navegador ``não mudou'', quase sempre\\
    o servidor antigo ainda está rodando.
  \end{alertblock}
\end{frame}

% =============================================================================
\begin{frame}{Bloco A --- Aquecimento}
  Ainda sem parâmetro --- só novos endereços.

  \vspace{0.4em}
  \begin{enumerate}
    \item \textbf{\texttt{/ola}} --- altere a mensagem para:\\
          \texttt{Olá, turma Entra21! Spring na veia.}
    \item \textbf{\texttt{/turma}} --- devolva o nome da disciplina\\
          e o nome do polo/cidade (texto livre).
    \item \textbf{\texttt{/hora}} --- devolva uma frase fixa, por exemplo:\\
          \texttt{A aula vai até as 22h.}
  \end{enumerate}

  \vspace{0.6em}
  Objetivo: criar \texttt{@GetMapping} novos com segurança.
\end{frame}

% =============================================================================
\begin{frame}[fragile]{Ferramenta nova: @RequestParam}
  Para ler valores da URL assim:

  \begin{center}
    \texttt{http://localhost:8080/saudacao?nome=Ana}
  \end{center}

\begin{lstlisting}[language=Java,basicstyle=\ttfamily\footnotesize]
@GetMapping("/saudacao")
public String saudacao(@RequestParam String nome) {
  return "Ola, " + nome + "!";
}
\end{lstlisting}

  \vspace{0.3em}
  Não esqueça do import:

\begin{lstlisting}[basicstyle=\ttfamily\footnotesize]
import org.springframework.web.bind.annotation.RequestParam;
\end{lstlisting}
\end{frame}

% =============================================================================
\begin{frame}{Bloco B --- Parâmetros na URL (?)}
  \begin{enumerate}
    \setcounter{enumi}{3}
    \item \textbf{\texttt{/saudacao?nome=}}\\
          Ex.: \texttt{?nome=Ana} $\rightarrow$ \texttt{Olá, Ana!}
    \item \textbf{\texttt{/dobro?numero=}}\\
          Ex.: \texttt{?numero=4} $\rightarrow$ \texttt{O dobro de 4 é 8.}\\
          Use \texttt{int} no \texttt{@RequestParam}.
    \item \textbf{\texttt{/soma?a=\&b=}}\\
          Ex.: \texttt{?a=2\&b=3} $\rightarrow$ \texttt{2 + 3 = 5}
    \item \textbf{\texttt{/media?a=\&b=}} (extra)\\
          Devolva a média dos dois números.
  \end{enumerate}
\end{frame}

% =============================================================================
\begin{frame}{Bloco C --- Decisão (if)}
  Java que vocês já sabem --- agora dentro do endpoint.

  \vspace{0.4em}
  \begin{enumerate}
    \setcounter{enumi}{7}
    \item \textbf{\texttt{/par-impar?numero=}}\\
          Par $\rightarrow$ \texttt{4 é par.} \quad
          Ímpar $\rightarrow$ \texttt{5 é ímpar.}
    \item \textbf{\texttt{/acesso?idade=}}\\
          \texttt{idade >= 18} $\rightarrow$ \texttt{Acesso liberado.}\\
          senão $\rightarrow$ \texttt{Acesso negado.}
    \item \textbf{\texttt{/desconto?valor=}}\\
          Se \texttt{valor >= 100}: 10\% de desconto e mostre o valor final.\\
          Senão: \texttt{Sem desconto. Valor: ...}
  \end{enumerate}
\end{frame}

% =============================================================================
\begin{frame}[fragile]{Bloco D --- Parâmetro opcional}
  Às vezes a pessoa acessa \textbf{sem} enviar o parâmetro.

\begin{lstlisting}[language=Java,basicstyle=\ttfamily\footnotesize]
@RequestParam(required = false) String nome
\end{lstlisting}

  \vspace{0.4em}
  \begin{enumerate}
    \setcounter{enumi}{10}
    \item \textbf{\texttt{/bemvindo}}
      \begin{itemize}
        \item Com \texttt{?nome=Joao} $\rightarrow$
              \texttt{Bem-vindo(a), Joao!}
        \item Sem nome $\rightarrow$
              \texttt{Bem-vindo(a), visitante!}
      \end{itemize}
  \end{enumerate}

  \vspace{0.5em}
  Este exercício é o ``ensaio'' do desafio de casa.
\end{frame}

% =============================================================================
\begin{frame}{Bloco E --- Se sobrar tempo (nível 2)}
  Para quem terminou os anteriores:

  \vspace{0.4em}
  \begin{enumerate}
    \setcounter{enumi}{11}
    \item \textbf{\texttt{/calc?a=\&b=\&op=}}\\
          \texttt{op} = \texttt{soma}, \texttt{sub}, \texttt{mult} ou \texttt{div}.\\
          Se dividir por zero: mensagem amigável (não deixe quebrar).
    \item \textbf{\texttt{/eco?texto=\&vezes=}}\\
          Repita o texto N vezes (uma frase por linha, usando \texttt{\textbackslash n}).
  \end{enumerate}

  \vspace{0.6em}
  \begin{block}{Meta da prática}
    Pedido na URL $\rightarrow$ Java processa $\rightarrow$ texto devolve no navegador.
  \end{block}
\end{frame}

% =============================================================================
\begin{frame}{Checklist rápido (se travar)}
  \begin{itemize}
    \item [ ] O método está dentro de uma classe com \texttt{@RestController}?
    \item [ ] O caminho do \texttt{@GetMapping} está igual ao da URL?
    \item [ ] Importei \texttt{RequestParam}?
    \item [ ] Reiniciei o Spring depois de salvar?
    \item [ ] Li o erro vermelho do terminal (de baixo para cima)?
  \end{itemize}

  \vspace{0.8em}
  Dúvida? Chame o professor --- mostre a URL e a tela do código.
\end{frame}

% #############################################################################
\section{Desafio (sozinhos)}

% =============================================================================
\begin{frame}{Desafio 01 --- Saudação personalizada}
  Criem um projeto (pode ser novo no Initializr) com:

  \vspace{0.5em}
  \begin{block}{Endpoint}
    \texttt{GET /saudacao?nome=Maria}
  \end{block}

  \vspace{0.3em}
  Resposta esperada:

  \begin{center}
    \texttt{Olá, Maria! Seja bem-vindo(a) ao Spring.}
  \end{center}

  \vspace{0.5em}
  Se abrir só \texttt{/saudacao} (sem nome):

  \begin{center}
    \texttt{Olá, visitante! Seja bem-vindo(a) ao Spring.}
  \end{center}
\end{frame}

% =============================================================================
\begin{frame}[fragile]{Desafio --- lembrete do @RequestParam}
  É o mesmo ideia do exercício \texttt{/bemvindo}, em projeto próprio:

\begin{lstlisting}[language=Java,basicstyle=\ttfamily\footnotesize]
@GetMapping("/saudacao")
public String saudacao(
    @RequestParam(required = false) String nome) {

  if (nome == null || nome.isBlank()) {
    nome = "visitante";
  }
  return "Ola, " + nome + "! Seja bem-vindo(a) ao Spring.";
}
\end{lstlisting}

  \vspace{0.3em}
  \begin{itemize}
    \item \texttt{required = false} permite acessar sem o parâmetro
    \item Faça \textbf{sozinho} (pode consultar os exercícios da aula)
  \end{itemize}
\end{frame}

% =============================================================================
\begin{frame}{Critérios de aceite do desafio}
  \begin{itemize}
    \item [ ] Projeto sobe com \texttt{.\textbackslash{}mvnw.cmd spring-boot:run}
    \item [ ] \texttt{/saudacao?nome=Ana} mostra Ana
    \item [ ] \texttt{/saudacao} mostra visitante
    \item [ ] Vocês explicam, em uma frase, o que \texttt{@GetMapping} faz
  \end{itemize}

  \vspace{0.8em}
  \begin{block}{Extra (opcional)}
    Criar \texttt{/soma?a=2\&b=3} retornando \texttt{5}.
  \end{block}
\end{frame}

% =============================================================================
\begin{frame}{Onde está o material}
  Tudo no repositório da disciplina:

  \vspace{0.5em}
  \begin{itemize}
    \item \texttt{slides/aula-01-introducao/} --- estes slides
    \item \texttt{ola-spring/} --- projeto feito juntos (consulta)
    \item \texttt{desafio-saudacao/} --- enunciado do desafio
  \end{itemize}

  \vspace{0.6em}
  Usem o projeto da aula como \textbf{gabarito vivo}, não como cópia cega:\\
  leiam, entendam, adaptem.
\end{frame}

% #############################################################################
\section{Recapitulando}

% =============================================================================
\begin{frame}{O que aprendemos hoje}
  \begin{enumerate}
    \item Spring Boot acelera o que já vimos em Java Web
    \item No Windows: JDK 21 + VS Code + extensões Java/Spring
    \item O Initializr gera o esqueleto do projeto
    \item \texttt{@RestController} + \texttt{@GetMapping} respondem no navegador
    \item \texttt{@RequestParam} lê valores da URL (\texttt{?nome=})
    \item Praticamos vários endpoints no \texttt{ola-spring}
  \end{enumerate}
\end{frame}

% =============================================================================
\begin{frame}{Para a próxima aula}
  \begin{itemize}
    \item Tragam o JDK e o VS Code já funcionando
    \item Se puderem, terminem o desafio da saudação
    \item Na próxima: vamos aprofundar Controllers e organização do código
  \end{itemize}

  \vspace{1em}
  \begin{center}
    \Large Obrigado --- e até a próxima!
  \end{center}
\end{frame}

\end{document}
